%%%%%%%%%%%%%%%%%%%%%%%%%%%%%%%%%%%%%%%%%%%%%%%%%%%%%
%			HLAVIČKA								%
%%%%%%%%%%%%%%%%%%%%%%%%%%%%%%%%%%%%%%%%%%%%%%%%%%%%%
\documentclass[openany,12pt]{memoir}

\usepackage[utf8]{inputenc} 
\usepackage[czech]{babel}
\usepackage[T1]{fontenc}
\usepackage[top=1.5cm, bottom=2cm, left=2cm, right=2cm]{geometry}  % --> NASTAVENÍ OKRAJŮ
\usepackage{fancyhdr}
\usepackage{graphicx}
\usepackage{lmodern}
\usepackage{xwatermark}
\usepackage{xcolor}
\usepackage{changepage}
\usepackage{newunicodechar}



%%%%%% Package na zpěvník
\usepackage[full]{leadsheets}%http://mirrors.nic.cz/tex-archive/macros/latex/contrib/leadsheets/leadsheets_en.pdf   --> dokumentace	
\definesongtitletemplate{empty}{} 
\setchords{
format = \bfseries,   %tučné akordy
minor = {mi},% 
input-notation = {german},%
output-notation = {german}%
}
\definesongtitletemplate{empty}{} 

\newlength{\drop}
\newwatermark[allpages,color=red!50,angle=0,scale=2, xpos=0,ypos=0]{\includegraphics[width=5cm]{obr/pozadi.jpg}} %--> dvojka na pozadí


%%%% Vlastní příkazy
\newcounter{Slokočet}   %Automatické číslování slok
\newcommand{\mezera}{\vspace*{0.5cm}}   %Horizontální odsazení slok
\newcommand{\stred}{5.2cm}   %%% Na zarovnání slok doprostřed, pozn. automatičtější zarovnávání na střed nejde
\newcommand{\refren}{\mezera \noindent \textbf{R:} } %refrén
\newcommand{\sloka}{\mezera \noindent \addtocounter{Slokočet}{1} \arabic{Slokočet}. } 	%sloka, která se automaticky čísluje
\newcommand{\ssloka}{\mezera \noindent}  % vlastní číslo sloky
\newcommand{\carka}{,\:}
\newcommand{\m}[1]{\color{white}{#1}}  %Pro akordy

%%%%%%%%%%%%%%%%%%%%%%%%%%%%%%%%%%%%%%%%%%%%%%%%%%%%%%
%			NÁVOD									 %
%%%%%%%%%%%%%%%%%%%%%%%%%%%%%%%%%%%%%%%%%%%%%%%%%%%%%%
%1. Věci v hlavičce IGNOROVAT
%2. Píseň psát do prostoru mezi \begin{song} a \end{song}
%3. další řádek se značí dvěma odsazeníma (= dvakrát stisknout enter)
%4. \refren vždy na začátku refrenu a \sloka na začátek sloky (automaticky se čísluje)
%  \ = alt gr + q ; [/] = alt gr f/g ; {/} = alt gr + b/n; ^ = alt gr + 3 + mezera
%Cokoliv napíšete do ^{  } se bude brát jako akord
%když se toto bude dotýkat nějakého slova (nebude mezi tím a slovem mezera)
%tak se akordy zjeví nad slovem, ale pište to před slovo
%Když se to nedotýká slova, tak akord lítá ve vzduchu a vytiskne se větší mezera
%První možnost je asi preferovanější
%5. Akordy stačí psát jen do první sloky, když se nezmění -- kytaristi to zvládnou



%<++>
%\usepackage{subfiles}
%</++>

\begin{document}

\pagestyle{empty}
%\addtocontents{toc}{\protect\thispagestyle{empty}}
%\tableofcontents \thispagestyle{empty}\newpage
\newgeometry{top=1.5cm, bottom = 0cm, left = 2cm, right = 2cm}




%Pro kompilaci je důležitý, aby ten input a \end{document} zůstali na posledních
%dvou řádcích!!!!
<<<<<<< HEAD
\begin{song}{title=\centering Muzeum \\\normalsize Jaromír Nohavica  \vspace*{-0.3cm}}  %% sem se napíše jméno songu a autor
\moveright 3cm \vbox{      %Varianta č. 1  ---> Jeden sloupec zarovnaný na střed	

\sloka
	^{D}Ve Slezském muzeu, ^{A}v oddělení ^{Hmi}třetihor
	
	je bílý ^{G}krokodýl a ^{Dmi}edvěd a liška a ^{A7}kamenní ^{D}trilobiti, ^{A7}
	
	^{D}chodí se tam jen tak, co ^{A}noha nohu ^{Hmi}mine, 
	
	abys viděl, jak ten ^{G}život plyne, ^{D}jaké je to všechno ^{A7}pomíjivé ^{D}živobytí,
	
	^{G}pak vyjdeš do parku a ^{D4sus}celou noc se ^{G}touláš noční Opavou

	a ^{C}opájíš se ^{G}představou, jaké to bude ^{D}v ^{G}ráji.
	

\refren
	V ^{D}pět třicet pět jednou z ^{A}pravidelných ^{Hmi}linek,
	
	sedm zastávek do ^{G}Kateřinek, ^{D}ukončete nástup, ^{A7}dveře se ^{D}zavírají.


\sloka
	Budeš-li poslouchat a nebudeš-li odmlouvat,
	
	složíš-li svoje maturity, vychováš pár dětí a vyděláš dost peněz,
	
	můžeš se za odměnu svézt na velkém kolotoči,
	
	dostaneš krásnou knihu s věnováním zaručeně,
	
	a ty bys chtěl plout na hřbetě krokodýla po řece Nil
	
	a volat: \uv{Tutanchámon, vivat, vivat!} po egyptském kraji.
	
\refren

\sloka
	Pionýrský šátek uvážeš si kolem krku,
	
	ve Valtické poručíš si čtyři deci rumu a utopence k tomu,
	
	na politém stole na ubruse píšeš svou rýmovanou Odysseu,
	
	nežli přijde někdo, abys šel už domů,
	
	ale není žádné doma jako není žádné venku, není žádné venku,
	
	to jsou jenom slova, která obrátit se dle libosti dají.

\refren

\sloka
	Možná si k tobě někdo přisedne,
	
	a možná to bude zrovna muž, který osobně znal Egypťana Sinuheta,
	
	dřevěnou nohou bude vyťukávat do podlahy rytmus metronomu,
	
	který tady klepe od počátku světa,
	
	nebyli jsme, nebudem a nebyli jsme, nebudem a co budem, až nebudem,
	
	jen navezená mrva v boží stáji.
	
	
\refren

}
\setcounter{Slokočet}{0}
\end{song}


=======
\begin{song}{title=\predtitle \centering Amerika \\\large Lucie  \vspace*{-0.3cm}}  %% sem se napíše jméno songu a autor
\begin{centerjustified}
\nejvetsi
\sloka 
	^{G\z }Nandej mi ^{D}do hlavy tvý ^{Ami\z }brouky 
	
	a Bůh nám seber ^{\z G}beznaděj.
	
	V duši zbylo ^{D\z }světlo z jedný ^{Ami}holky,
	
	tak mi teď za to ^{\z G}vynadej.
	
	Zima a ^{D\z }promarněný ^{Ami}touhy, 
	
	do vrásek stromů padá ^{G}déšť. 
	
	Zbejvaj roky ^{D}asi ne moc ^{Ami}dlouhý.
	
	Do vlasů mi ^{\z C}zabroukej pá pa pá ^{G}pá. 

\refren
/: Pá ^{G/F#}pá pa ^{Emi\z}pá,~pá pa pa ^{G}pá. :/

\sloka
	Tvoje oči jenom žhavý tóny, 
	
	dotek slunce zapadá.
	
	Horkej vítr rozezní mý zvony.
	
	Do vlasú ti zabrouká pá pa pa pá. 

\refren

\sloka
	Na obloze křídla tažnejch ptáků,
	
	tak už na svý bráchy zavolej. 
	
	Na tváře ti padaj slzy z mraků
	
	a Bůh nám sebral beznaděj. 
	
	V duši zbylo světlo z jedný holky,
	
	do vrásek stromů padá déšť. 
	
	Poslední dny hodiny a roky, 
	
	do vlasů ti zabrouká pá pa pa pá.

\refren


\end{centerjustified}
\setcounter{Slokočet}{0}
\end{song}



>>>>>>> e6ebb57d3685c545b7aef13aad0b76371c34b457
\end{document}
